\subsection{Pregunta y propósito analítico}

El proyecto busca describir cómo se distribuyen territorialmente las condiciones socioeconómicas observadas en 2018 y, en una fase posterior, estudiar su asociación con la contratación pública registrada en SECOP II durante 2026. La vulnerabilidad se entiende como una condición multidimensional relacionada con recursos, capacidades y exposición a restricciones; no es sinónimo automático de pobreza monetaria ni una característica individual que pueda inferirse directamente de un promedio municipal \citep{sen1999development,kaztman1999activos,pizarro2001vulnerabilidad}. El IPM municipal se conserva como referente externo para contraste posterior, no como insumo del PCA, para evitar validar una medida con componentes que reproducen su definición.

El estado del arte del proyecto aporta antecedentes de desigualdad regional y análisis territorial de pobreza y contratación. Se usan como contexto para formular la pregunta y justificar una lectura espacial, no como evidencia de que una combinación particular de variables sea válida para este IVS \citep{bonet2006desequilibrios,monsalvo2025ipm}.

\subsection{Diseño y unidad de inferencia}

La fase censal es transversal, descriptiva y ecológica: cada fila representa un código territorial y las observaciones describen agregados del CNPV 2018. Las asociaciones obtenidas no permiten atribuir características a hogares o personas específicas ni sostener causalidad. El IVS pretende resumir diferencias territoriales, no medir directamente la vulnerabilidad de cada residente.

El corte 2018 es la referencia socioeconómica del índice. SECOP II se mantendrá como bloque de análisis posterior y, por decisión del equipo, se priorizarán contratos de 2026. La diferencia temporal entre 2018 y 2026 exige interpretar el resultado como asociación entre una condición censal de referencia y una distribución contractual posterior, no como comparación contemporánea ni efecto causal.

\subsection{Fundamento estadístico de un índice compuesto}

Un indicador compuesto requiere definir constructo, cobertura, selección, normalización, agregación y sensibilidad; una técnica de ponderación no reemplaza las decisiones conceptuales \citep{booysen2002overview,oecd2008composite}. El PCA es una transformación lineal que produce combinaciones ortogonales ordenadas por varianza explicada. No identifica por sí mismo dimensiones de vulnerabilidad ni asigna pesos normativos justos. Por ello, se propone como procedimiento exploratorio condicionado a la coherencia de las variables, su estructura de correlación, su estabilidad y la interpretación de las cargas \citep{jolliffe2002principal,abdi2010principal}.

La escala 0--1 facilita leer cada proporción y orientar su significado. La matriz z pone las variables en varianza unitaria para que las diferencias de unidades no determinen la solución. La estandarización no homogeneiza la calidad del dato, no convierte una proxy en constructo validado y no significa que el efecto sustantivo de todas las variables deba ser igual.

\subsection{Trazabilidad de estados de evidencia}

En el informe se distinguen tres niveles: \textbf{verificado en archivo} (se reproduce desde una fuente o salida), \textbf{decisión acordada} (regla analítica adoptada por el equipo) e \textbf{interpretación o hipótesis} (lectura conceptual que deberá contrastarse). Una comprobación aritmética no se presenta como validación teórica; una decisión de ponderación tampoco se atribuye al DANE.
